\chapter{Complexity — Search}
%%% generated by the blueprint pipeline from TCSlib.Complexity.SpaceComplexity.ExpBFS.Search
\section{Overview}

Declarations of the Lean module \texttt{TCSlib.Complexity.SpaceComplexity.ExpBFS.Search}.

\section{Declarations}

\begin{definition}[Number of Unmarked Codes]
\lean{Complexity.LogProg.ExpBFS.unmarked}
\label{Complexity.LogProg.ExpBFS.unmarked}
\leanok
\uses{Complexity.LogProg.ExpBFS.cell}
For natural numbers $b,K$ and a Boolean-valued memory function $M:\mathbb
Z\to\{\mathrm{false},\mathrm{true}\}$, define the number of unmarked codes to be the
cardinality of the set of integers $v$ with $0\le v<K$ for which
$M(\mathrm{cell}(b,v))=\mathrm{false}$.
\end{definition}

\begin{lemma}[Strict decrease in the unmarked count]
\lean{Complexity.LogProg.ExpBFS.unmarked_lt}
\label{Complexity.LogProg.ExpBFS.unmarked_lt}
\leanok
\uses{Complexity.LogProg.ExpBFS.Mono, Complexity.LogProg.ExpBFS.Out, Complexity.LogProg.ExpBFS.cell, Complexity.LogProg.ExpBFS.unmarked}
\difficulty{5}
Let $b,K$ be natural numbers and let $M,M':\mathbb Z\to\{\mathrm{false},\mathrm{true}\}$
be markings of integer cells. Suppose every cell marked by $M$ is also marked by $M'$,
and the two markings agree at every cell other than the cells $b,b+1,\ldots,b+K-1$. If
$M'\ne M$, then the number of unmarked cells in this finite range is strictly smaller
for $M'$ than for $M$.
\end{lemma}

\begin{theorem}[Closure after a relaxation round]
\lean{Complexity.LogProg.ExpBFS.rounds_run}
\label{Complexity.LogProg.ExpBFS.rounds_run}
\leanok
\uses{Complexity.LogProg.ExpBFS.A, Complexity.LogProg.ExpBFS.G, Complexity.LogProg.ExpBFS.Mono, Complexity.LogProg.ExpBFS.Mono.refl, Complexity.LogProg.ExpBFS.Mono.trans, Complexity.LogProg.ExpBFS.Out, Complexity.LogProg.ExpBFS.Out.refl, Complexity.LogProg.ExpBFS.Out.trans, Complexity.LogProg.ExpBFS.SBd, Complexity.LogProg.ExpBFS.St, Complexity.LogProg.ExpBFS.St.Bd, Complexity.LogProg.ExpBFS.cell, Complexity.LogProg.ExpBFS.cf, Complexity.LogProg.ExpBFS.cont, Complexity.LogProg.ExpBFS.ex1_pre, Complexity.LogProg.ExpBFS.outer_run, Complexity.LogProg.ExpBFS.reach_cf, Complexity.LogProg.ExpBFS.rr, Complexity.LogProg.ExpBFS.step_of, Complexity.LogProg.ExpBFS.stp, Complexity.LogProg.ExpBFS.unmarked, Complexity.LogProg.ExpBFS.unmarked_lt, Complexity.LogProg.MReach, Complexity.LogProg.MReach.refl, Complexity.LogProg.astep}
\difficulty{6}
Fix a natural parameter $K_0$, an oracle, a Boolean string $x$, and natural numbers
$B,P$. Let $\varphi\in\{\mathrm{false},\mathrm{true}\}$, let $j$ be a natural number,
let $S$ be a machine state, and let $M:\mathbb{Z}\to\{\mathrm{false},\mathrm{true}\}$
satisfy that exactly $j$ of the $K$ cells at positions $B_s+v$ for $v<K$ are unmarked.
Suppose $U=K$, $V=T=T_2=Z=0$, $ch\leq 1$, every component of $S$ is at most $B$, and
$B_s+K\leq B$, $B_s+K\leq P$, and $K+1\leq B$. Then there is a memory $M'$ such that,
starting at label $\mathrm{rl0}(\varphi)$ with state $S$, memory $M$, and head position
$0$, the machine reaches label $\mathrm{cn0}(\varphi)$ with state $S$ except that $ch$
is reset to $0$, memory $M'$, and head position $0$, while its guard holds at every time
strictly before arrival. Moreover, $M'$ extends the marked positions of $M$ and agrees
with $M$ outside the block of positions $B_s+v$ with $v<K$. For every predicate $Q$ on
natural numbers, if $Q$ is preserved by each true step between indices below $K$ and
holds at every initially marked index below $K$, then it holds at every index below $K$
marked in $M'$. Also, for all $a,b<K$, if $a$ is marked in $M'$ and the step from $a$ to
$b$ is true, then $b$ is marked in $M'$.
\end{theorem}

\begin{lemma}[Filtered interval cardinality recurrence]
\lean{Complexity.LogProg.ExpBFS.card_filter_Ico_succ}
\label{Complexity.LogProg.ExpBFS.card_filter_Ico_succ}
\leanok
\difficulty{3}
Let $p$ be a decidable predicate on the natural numbers, and let $u<K$ be natural
numbers. The number of elements of the half-open interval $[u,K)$ satisfying $p$ equals
$1$ if $p(u)$ holds and $0$ otherwise, plus the number of elements of $[u+1,K)$
satisfying $p$.
\end{lemma}

\begin{theorem}[Marked-Code Counting Run]
\lean{Complexity.LogProg.ExpBFS.count_run}
\label{Complexity.LogProg.ExpBFS.count_run}
\leanok
\uses{Complexity.LogProg.ExpBFS.A, Complexity.LogProg.ExpBFS.G, Complexity.LogProg.ExpBFS.SBd, Complexity.LogProg.ExpBFS.Site.idx, Complexity.LogProg.ExpBFS.Site.set, Complexity.LogProg.ExpBFS.St, Complexity.LogProg.ExpBFS.St.Bd, Complexity.LogProg.ExpBFS.St.toV_U, Complexity.LogProg.ExpBFS.acExit, Complexity.LogProg.ExpBFS.ac_run, Complexity.LogProg.ExpBFS.card_filter_Ico_succ, Complexity.LogProg.ExpBFS.cell, Complexity.LogProg.ExpBFS.cf, Complexity.LogProg.ExpBFS.reach_cf, Complexity.LogProg.ExpBFS.step_of, Complexity.LogProg.MReach, Complexity.LogProg.MReach.refl, Complexity.LogProg.MReach.trans, Complexity.LogProg.astep}
\difficulty{6}
Fix natural numbers $K_0,B,P$, an oracle
$\mathrm{orc}:\{0,1,2,3\}\times\{0,1\}^*\to\{\mathrm{false},\mathrm{true}\}$, and a
finite bit string $x$. For a Boolean $\varphi$, a natural number $j$, a register state
$S$, and a memory function $M:\mathbb Z\to\{\mathrm{false},\mathrm{true}\}$, suppose
that $S.K-S.U=j$ (with truncated natural-number subtraction), $S.U\le S.K$,
$S.T=S.T_2=S.Z=0$, and $S.c+(S.K-S.U)\le S.K$. Suppose also that every register of $S$
is at most $B$, that $S.B_s+S.K\le B$, $S.B_s+S.K\le P$, and $S.K+1\le B$. Then the
machine has a finite run from the count-search configuration with registers $S$, memory
$M$, and memory-head position $0$ to the branch for $\varphi$, with $U=S.K$ and $c$
equal to its initial value plus the number of integers $v$ in $[S.U,S.K)$ for which
$M(S.B_s+v)=\mathrm{true}$. Every other register retains its initial value, and the
safety guard holds at every configuration before the endpoint.
\end{theorem}

\begin{theorem}[Bounded reachability search]
\lean{Complexity.LogProg.ExpBFS.search_run}
\label{Complexity.LogProg.ExpBFS.search_run}
\leanok
\uses{Complexity.LogProg.ExpBFS.A, Complexity.LogProg.ExpBFS.Cl, Complexity.LogProg.ExpBFS.Cl.step, Complexity.LogProg.ExpBFS.G, Complexity.LogProg.ExpBFS.Mono, Complexity.LogProg.ExpBFS.Mono.trans, Complexity.LogProg.ExpBFS.Out, Complexity.LogProg.ExpBFS.Out.trans, Complexity.LogProg.ExpBFS.SBd, Complexity.LogProg.ExpBFS.St, Complexity.LogProg.ExpBFS.St.Bd, Complexity.LogProg.ExpBFS.cell, Complexity.LogProg.ExpBFS.cf, Complexity.LogProg.ExpBFS.count_run, Complexity.LogProg.ExpBFS.ini_run, Complexity.LogProg.ExpBFS.reach_cf, Complexity.LogProg.ExpBFS.rounds_run, Complexity.LogProg.ExpBFS.rr, Complexity.LogProg.ExpBFS.step_of, Complexity.LogProg.ExpBFS.stp, Complexity.LogProg.MReach, Complexity.LogProg.MReach.trans, Complexity.LogProg.astep}
\difficulty{6}
Let $K_0$ be a natural number, let $\mathrm{orc}$ assign a Boolean to each of four query
indices and each finite bit string, let $x$ be a finite bit string, and let $B,P$ be
natural numbers. Let $\varphi$ be a Boolean, let $S$ be a machine state, and let
$M:\mathbb{Z}\to\{\mathrm{false},\mathrm{true}\}$ be a memory function. Suppose
$S.U=S.V=S.T=S.T_2=S.Z=S.ch=S.c=0$, $S$ satisfies $\mathrm{SBd}(B,P)$, and
$M(S.Bs+v)=\mathrm{false}$ for every natural number $v<S.K$. Then there is a memory
function $M'$ such that the machine has a finite run from the initial configuration to
the configuration at the second count-comparison label if $\varphi$ is true, and at the
first-search cleanup label otherwise; every configuration before the endpoint satisfies
the guard $G(x,B,P)$. At the endpoint, $U=S.K$ and $c$ is the number of indices $v<S.K$
reachable from some $u<S.K$ with a true initial query by a finite sequence of true step
queries whose successive indices remain below $S.K$; all other registers retain their
values from $S$. Moreover, every position marked true by $M$ is marked true by $M'$, no
memory cell outside $\{S.Bs+v:v<S.K\}$ changes, and for each $v<S.K$, $M'(S.Bs+v)$ is
true exactly when $v$ has that reachability property.
\end{theorem}
